\documentclass[10pt]{article}
\usepackage[T1]{fontenc}
\usepackage[utf8]{inputenc}
\usepackage[french]{babel}
\usepackage{amsmath,amssymb,mathrsfs,mathpazo}
\usepackage{geometry}
\geometry{a4paper,margin=2cm}
\usepackage{graphicx,array,tabularx,tabularray,colortbl,multirow,multicol,booktabs,enumitem,pifont,ulem,fancybox,diagbox,cancel}
\usepackage{pstricks,pstricks-add,pst-plot,pst-tree,pst-node,pst-text,pst-func,pst-math,pst-bspline,pst-3dplot,pst-eucl,pst-solides3d}
\let\pscalculate\undefined
\usepackage{pst-bezier}
\usepackage{tikz,pgfplots,tkz-tab}
\usetikzlibrary{arrows,arrows.meta,patterns,calc,positioning,trees,decorations.markings,shapes.geometric,shapes.misc,shapes.arrows,intersections,decorations.pathmorphing,decorations.pathreplacing,backgrounds,shadows,babel}
\pgfplotsset{compat=1.18}
\usepgfplotslibrary{fillbetween}
\usepackage[np]{numprint}
\usepackage{listings,makecell,pas-tableur,twoopt}
\let\barre\undefined
\usepackage[pstricks]{tablvar}
\definecolor{darkgreen}{RGB}{0,100,0}
\providecommand{\red}{\color{red}}
\providecommand{\blue}{\color{blue}}
\providecommand{\green}{\color{green}}
\providecommand{\black}{\color{black}}
\providecommand{\white}{\color{white}}
\providecommand{\gray}{\color{gray}}
\providecommand{\cyan}{\color{cyan}}
\providecommand{\magenta}{\color{magenta}}
\providecommand{\yellow}{\color{yellow}}
\providecommand{\violet}{\color{violet}}
\providecommand{\orange}{\color{orange}}
\providecommand{\euro}{\texteuro}
\providecommand{\EUR}{\texteuro}
\providecommand{\ieme}{\textsuperscript{e}}
\providecommand{\degres}{\ensuremath{^\circ}}
\providecommand{\degree}{\ensuremath{^\circ}}
\providecommand{\vv}[1]{\overrightarrow{#1}}
\providecommand{\Coord}[1]{#1}
\providecommand{\decosix}{}
\providecommand{\starredbullet}{\ensuremath{\bullet}}
\pagestyle{empty}
\setlength{\parindent}{0pt}
\usepackage{adjustbox}\DeclareUnicodeCharacter{2236}{:}\DeclareUnicodeCharacter{2190}{\ensuremath{\leftarrow}}
\providecommand{\R}{\mathbb{R}}
\providecommand{\N}{\mathbb{N}}
\providecommand{\Z}{\mathbb{Z}}
\providecommand{\Q}{\mathbb{Q}}
\providecommand{\D}{\mathbb{D}}
\providecommand{\C}{\mathbb{C}}
\providecommand{\e}{{\mathrm{e}}}
\providecommand{\vect}[1]{\overrightarrow{#1}}
\providecommand{\vectt}[1]{\overrightarrow{\mathrm{#1}}}
\providecommand{\vvt}[1]{\overrightarrow{\mathrm{#1}}}
\providecommand{\np}[2][]{{#2~#1}}
\providecommand{\vv}[1]{{\overrightarrow{#1}}}
\renewcommand{\d}{{\mathrm{d}}}
\providecommand{\degres}{{^\circ}}
\providecommand{\degre}{{^\circ}}
\providecommand{\blue}{\color{blue}}
\providecommand{\red}{\color{red}}
\providecommand{\enskip}{\quad}
\providecommand{\llbracket}{\lbrack\!\lbrack}
\providecommand{\rrbracket}{\rbrack\!\rbrack}
\providecommand{\euro}{{\text{€}}}
\providecommand{\no}{{\text{n°}}}
\providecommand{\interval}[2]{{[#1;#2]}}
\providecommand{\intervFO}[2]{{[#1;#2[}}
\providecommand{\dotfill}{\ldots}

\providecommand{\ds}{\displaystyle}
\providecommand{\textding}[1]{#1}
\providecommand{\Oij}{\(\mathrm{O};\vec{\imath},\vec{\jmath}\)}
\providecommand{\Oijk}{\(\mathrm{O};\vec{\imath},\vec{\jmath},\vec{k}\)}
\providecommand{\Ouv}{\(\mathrm{O};\vec{u},\vec{v}\)}
\providecommand{\intervFF}[2]{\ensuremath{[#1;#2]}}
\renewcommand{\euro}{\ifmmode\text{\texteuro}\else\texteuro\fi}
\DeclareUnicodeCharacter{00B0}{\ifmmode^{\circ}\else\textdegree\fi}
\DeclareUnicodeCharacter{20AC}{\ifmmode\text{\texteuro}\else\texteuro\fi}
% Énoncés : Éducation nationale. Transcription et corrigés : APMEP, Denis Vergès et contributeurs. Découpage et classement : CoopMaths, Sébastien Lozano, Éric Elter et contributeurs. Adaptation de la présentation : Maths Explorer.
% CC BY-SA : https://coopmaths.fr/www/about/
% https://coopmaths.fr/alea/static/eam/2025/tex/eam_technologique_2025_06_sujet0v2_1.tex
\begin{document}
\noindent\textit{APMEP / CoopMaths. CC BY-SA. Présentation adaptée par Maths Explorer.}\par\medskip
\textbf{Pour cette première partie, aucune justification n'est demandée et une seule réponse est possible par question. Pour chaque question, reportez son numéro sur votre copie et indiquez votre réponse.}

\subsubsection*{Question 1}
Un article coûte 400 euros. Le prix augmente de 20\%. Le nouveau prix est

\textbf{a.} 420 euros \hfill \textbf{b.} 480 euros \hfill \textbf{c.} 500 euros \hfill \textbf{d.} 320 euros

\subsubsection*{Question 2}
Un sac coûte 130 euros. Le prix baisse de 10\%. Le nouveau prix est

\textbf{a.} $130 \times 0,1$ \hfill \textbf{b.} $130 \times \left(-\dfrac{10}{100}\right)$ \hfill \textbf{c.} $130 \times \left(1 + \dfrac{10}{100}\right)$ \hfill \textbf{d.} $130 \times 0,9$

\subsubsection*{Question 3}
Le prix d'un article est noté $P$. Il connaît deux augmentations de 20\%. \\
Le prix après ces augmentations est

\textbf{a.} $P \times \left(1 + \left(\dfrac{20}{100}\right)^2\right)$ \hfill \textbf{b.} $P \times 1,40$ \hfill \textbf{c.} $\dfrac{P}{1,44}$ \hfill \textbf{d.} $P \times 1,2^2$

\subsubsection*{Question 4}
Lors d'une élection, le quart des électeurs a voté pour A, 20\% a voté pour B, un tiers a voté pour C, et le reste a voté pour D. \\
Le candidat ayant recueilli le \underline{moins} de votes est

\textbf{a.} A \hfill \textbf{b.} B \hfill \textbf{c.} C \hfill \textbf{d.} D

\subsubsection*{Question 5}
On considère $A = \dfrac{2}{1 - \dfrac{2}{3}}$. On a

\textbf{a.} $A = -1$ \hfill \textbf{b.} $A = \dfrac{2}{3}$ \hfill \textbf{c.} $A = 6$ \hfill \textbf{d.} $A = 9$

\subsubsection*{Question 6}
On considère $A = \dfrac{1}{100} + \dfrac{1}{1\,000}$. On a

\textbf{a.} $A = 100,001$ \hfill \textbf{b.} $A = \dfrac{2}{100\,000}$ \hfill \textbf{c.} $A = 0,11$ \hfill \textbf{d.} $A = 0,011$

\subsubsection*{Question 7}
Une durée de 75 minutes correspond à

\textbf{a.} 1,15 heure \hfill \textbf{b.} 1,25 heure \hfill \textbf{c.} 0,75 heure \hfill \textbf{d.} 1,4 heure

\subsubsection*{Question 8}
$10^{30} + 10^{-30}$ est environ égal à

\textbf{a.} $10^0$ \hfill \textbf{b.} 0 \hfill \textbf{c.} $10^{30}$ \hfill \textbf{d.} $20^{30}$

\subsubsection*{Question 9}
La seule droite pouvant correspondre à l'équation $y = -2x + 5$ est

\begin{minipage}[t]{0.24\textwidth}
\begin{center}
\textbf{a.} la droite $D_1$

\medskip

\begin{tikzpicture}[scale=0.8]
\draw[thick, ->] (-2,0) -- (2,0);
\draw[thick, ->] (0,-1) -- (0,2);
\draw[very thick, color=red] (-1,-1) -- (2,1);
\end{tikzpicture}
\end{center}
\end{minipage}
\hfill
\begin{minipage}[t]{0.24\textwidth}
\begin{center}
\textbf{b.} la droite $D_2$

\medskip

\begin{tikzpicture}[scale=0.8]
\draw[thick, ->] (-2,0) -- (2,0);
\draw[thick, ->] (0,-1) -- (0,2);
\draw[very thick, color=red] (-2,-0.5) -- (2,1.75);
\end{tikzpicture}
\end{center}
\end{minipage}
\hfill
\begin{minipage}[t]{0.24\textwidth}
\begin{center}
\textbf{c.} la droite $D_3$

\medskip

\begin{tikzpicture}[scale=0.8]
\draw[thick, ->] (-2,0) -- (2,0);
\draw[thick, ->] (0,-1) -- (0,2);
\draw[very thick, color=red] (-1.5,2) -- (2,-1);
\end{tikzpicture}
\end{center}
\end{minipage}
\hfill
\begin{minipage}[t]{0.24\textwidth}
\begin{center}
\textbf{d.} la droite $D_4$

\medskip

\begin{tikzpicture}[scale=0.8]
\draw[thick, ->] (-2,0) -- (2,0);
\draw[thick, ->] (0,-1) -- (0,2);
\draw[very thick, color=red] (-2,1) -- (0.5,-1);
\end{tikzpicture}
\end{center}
\end{minipage}

\subsubsection*{Question 10}
La solution de l'équation $3x = 0$ est

\textbf{a.} $x = -3$ \hfill \textbf{b.} $x = \dfrac{1}{3}$ \hfill \textbf{c.} $x = -\dfrac{1}{3}$ \hfill \textbf{d.} $x = 0$

\subsubsection*{Question 11}
La solution de l'équation $\dfrac{144}{x} = 9$ est

\textbf{a.} $x = 144 \times 9$ \hfill \textbf{b.} $x = \dfrac{9}{144}$ \hfill \textbf{c.} $x = \dfrac{144}{9}$ \hfill \textbf{d.} $x = -16$

\subsubsection*{Question 12}
Voici les notes sur vingt obtenues par un élève en mathématiques :

\begin{center}
\begin{tabular}{|c|c|c|c|c|}
\hline
Note & 10 & 13 & 12 & $x$ \\
\hline
Coefficient & 1 & 1 & 1 & 2 \\
\hline
\end{tabular}
\end{center}

On cherche ce que doit valoir $x$ pour que la moyenne de l'élève soit égale à 15.

\textbf{a.} $x = 20$ \hfill \textbf{b.} $x = 18$

\textbf{c.} $x = 15$ \hfill \textbf{d.} Impossible : il faudrait une note strictement supérieure à vingt.

\end{document}
